\subsection{Taps, reverse feed, autotransformers and standards}
\label{sec:na-application}

\subsubsection{Taps}
The 480\,V winding of a general-purpose dry-type transformer has two taps
of 2.5\,\% above and two below the rated voltage; some designs have two
above and four below~[7]. The taps are set off-circuit.

\subsubsection{Reverse feed}
A unit energised from its 240\,V winding to supply 480\,V operates as a
step-up transformer. The inrush current is higher than in the rated
direction, and the taps then act on the output voltage; reverse feed
therefore needs the manufacturer's approval or a unit designed for it.

\subsubsection{Autotransformer}
Where isolation is not required, an autotransformer changes the voltage
with a winding power of
\begin{equation}
  S_{w} = S_{load} \cdot \left(1 - \frac{U_{2}}{U_{1}}\right)
  \label{eq:na-auto}
\end{equation}
with $U_{1}$ the higher and $U_{2}$ the lower voltage. For 480\,V to
240\,V, $S_{w} = 0.5 \cdot S_{load}$; for a buck-boost connection with a
small voltage change, $S_{w}$ is a few per cent of the load. The secondary is
not a separately derived system, and the restrictions on autotransformers
of~[3, 210.9 and 450.4] apply.

\subsubsection{Standards}
General-purpose dry-type transformers are built to IEEE~C57.12.01~[8]
and listed to UL~1561~[9]. Their minimum efficiency is set by
10~CFR~431~[10]. For non-linear loads the rating is specified as a
K-factor, and the derating follows IEEE~C57.110~[11].
\FloatBarrier
